\section{Random digits, complex heights and echoes}
\label{play:section}

The same constructions behave differently under three familiar tests.
Binary normality asks for every finite word; complex value distribution
asks which heights a curve reaches; Fourier decay asks whether a
measure forgets high frequencies.  The next observations connect those
tests to arithmetic already built into our objects.

\subsection{What would genuinely random digits mean?}

The original constant began by repeating Goldbach answers.  If every
answer is yes, its digits become $(1-\lambda(n))/2$.  It is tempting
to call these digits random.  Ordinary normality gives that word a
precise meaning: every binary word of length $k$ must occur with
frequency $2^{-k}$, for every $k$.

The relevant classical assertion is the ordinary, naturally averaged
Chowla conjecture for the Liouville function.  In its sign-pattern
form it predicts precisely these frequencies; see
\cite[Conjecture 1.3]{playMatomakiRadziwillTao2015}.  Its correlation
form says that for every nonempty finite set $J$ of distinct
nonnegative integers,
\begin{equation}
 \frac1N\sum_{n\le N}\prod_{j\in J}\lambda(n+j)
 \longrightarrow0.
 \label{play:chowla}
\end{equation}
In particular, this is a fixed-step assertion.  Averaging over the
step in the sampling theorem does not establish it.

\begin{proposition}\label{play:normality}
The original Goldbach constant is normal in base two if and only if
Goldbach's conjecture and the Liouville Chowla conjecture
\eqref{play:chowla} both hold.
\end{proposition}
\begin{proof}
Normality implies simple normality, which already implies Goldbach
by Theorem~\ref{fnd:balance}.  Under Goldbach the digit at $n$ is
$(1-\lambda(n))/2$.  For a word $w=(w_0,\ldots,w_{k-1})$, its
indicator at $n$ is
\[
 2^{-k}\prod_{j=0}^{k-1}
       \bigl(1+(-1)^{w_j}\lambda(n+j)\bigr).
\]
Expanding this product shows that \eqref{play:chowla} gives the
frequency $2^{-k}$.  Conversely, uniform frequencies of all words
give zero average to every nonconstant product of their signs.
Every finite $J$ lies inside a consecutive block, so this recovers
\eqref{play:chowla}.
\end{proof}

This equivalence helps locate the difficulty.  Normality of this one
real number simultaneously asks that no Goldbach answer be missing
and that every finite pattern of Liouville signs have its expected
frequency.  It is a meeting point of two existing conjectures, rather
than a newly available shortcut to either one.  The primitive
sampling theorem supplies a useful contrast: its fair-bit law is
unconditional because it averages over an additional variable.


\subsection{Every complex height returns}

On the real line our prime density is a positive decreasing curve.
Away from that line, its entire continuation behaves very differently.
Put
\[
 m=f(0)=\frac1{2\mathfrak C_2}.
\]
The reflection law $f(z)+f(-z)=2m$ has a surprisingly strong
consequence when combined with Picard's theorem.

\begin{proposition}\label{play:picard}
For every $a\in\mathbb C$ other than possibly $m$, the equation
$f(z)=a$ has infinitely many solutions.  In particular $f$ has
infinitely many complex zeros, but none on the real or imaginary
axis.
\end{proposition}
\begin{proof}
The nonzero coefficients at every prime position show that $f$ is
transcendental entire.  Picard's theorem says that such a function
can have at most one complex value with only finitely many
preimages \cite[Chapter II, Section 9]{playPicard1880}.  Reflection
pairs the solutions of $f(z)=a$ with those of $f(z)=2m-a$.
If $a$ had only finitely many preimages, so would $2m-a$;
hence these values must coincide, giving $a=m$.

The real-axis limits and monotonicity in
Corollary~\ref{fnd:reflection} imply $0\le f(x)\le2m$ there.
Equality at either endpoint would make $f$ constant on a real
half-line, hence constant everywhere by the identity theorem.
Thus $0<f(x)<2m$ for every real $x$.  Finally the Taylor
coefficients are real, so
$f(-iy)=\overline{f(iy)}$.  Reflection therefore gives
$\Re f(iy)=m>0$, excluding imaginary-axis zeros as well.
\end{proof}

The conclusion uses a classical theorem, but it sharpens the picture
of this particular density: a curve positive everywhere on the real
line has infinitely many zeros off both coordinate axes.  Exactly
one height escapes the argument, its central height $m$.

Reflection alone cannot settle the middle height: an odd entire
function can have only one zero.  We initially approached it as a
question about how the prime gaps constrain an entire function.
But a simpler feature of the primes settles it, and more: the only
prime divisible by $3$ is $3$.  Among the exponents $p-2$, therefore,
only the exponent $1$ is congruent to $1$ modulo $3$.  With
$\omega=e^{2\pi i/3}$, rotating the argument through three equally
spaced angles and taking a weighted sum gives
\[
 f(z)+\omega^2 f(\omega z)+\omega f(\omega^2z)
   =\frac{3z}{K'(-2)}.
\]
Every term except the linear one cancels.  This finite identity
contains information that the size of the prime gaps alone misses.
Using larger primes in the same way gives the following conclusion.

\begin{theorem}[Every polynomial target is met infinitely often]
\label{play:midpoint}
For every polynomial $P\in\mathbb C[z]$, the equation
\[
 f(z)=P(z)
\]
has infinitely many distinct solutions in $\mathbb C$.
In particular every complex value, including the middle height
$m=f(0)$, is attained infinitely often.
\end{theorem}
\begin{proof}
Choose an odd prime $q>\deg P+2$; when $P=0$, any odd prime will
do.  Put $\zeta=e^{2\pi i/q}$ and $h=f-P$.
The root-of-unity filter gives
\begin{equation}
 \sum_{j=0}^{q-1}\zeta^{-j(q-2)}h(\zeta^jz)
    =\frac{qz^{q-2}}{K'(1-q)}.
 \label{play:prime-filter}
\end{equation}
Indeed, the weighted sum kills a monomial $z^n$ unless
$n\equiv q-2\pmod q$.  It kills all of $P$, while a surviving
exponent $p-2$ of $f$ would require $q\mid p$, hence $p=q$.
The coefficient on the right is nonzero because the zero of $K$
at $1-q$ is simple.

Suppose $h$ had finitely many zeros.  Factoring them out, with
multiplicity, gives $h=D e^H$, where $D$ is a nonzero polynomial
and $H$ is entire: the remaining zero-free entire function has an
entire logarithm.  The function $H$ is nonconstant, since $h$ is
transcendental.  Substitution into \eqref{play:prime-filter} gives
a polynomial-coefficient relation among
\[
 e^{H(z)},\ e^{H(\zeta z)},\ldots,
 e^{H(\zeta^{q-1}z)},\ e^0.
\]
Narasimhan's theorem says that exponentials of entire functions
whose pairwise differences are nonconstant are linearly independent
over $\mathbb C[z]$ \cite[Theorem 1]{playNarasimhan1971}.
To apply it, group the rotated exponents whenever their difference
is constant, absorbing the corresponding constant exponential
factors into their polynomial coefficients.  Discard any group
whose combined coefficient vanishes.  The exponent $0$ remains
in its own group, since every $H(\zeta^jz)$ is nonconstant.  Its
coefficient is $-qz^{q-2}/K'(1-q)$, which is nonzero.
The resulting relation contradicts Narasimhan's theorem.
\end{proof}

Thus the bounded real curve does more than revisit every complex
height.  It meets every prescribed polynomial target infinitely
often.  The proof uses a classical independence theorem; its
arithmetic input is the elementary fact that the only prime divisible
by a prime $q$ is $q$ itself. No estimate for the distribution of primes is
needed in this step.

There is a small geometric consequence of the same cancellation.
Every complex height occurs infinitely often, yet it cannot occur
at all three vertices of an equilateral triangle centred at the
origin.  More generally, no nondegenerate regular polygon with an
odd prime number $q$ of vertices, centred at the origin, has all its
vertices at one height of $f$.  Indeed, if $f(\zeta^jz)=a$ for
every $j$ and $z\ne0$, the left side of
\eqref{play:prime-filter}, with $P=a$, is zero and the right side
is not.

The infinitude conclusion also survives any number of derivatives:
for every integer $r\ge0$ and every polynomial $P$, the equation
$f^{(r)}(z)=P(z)$ has infinitely many solutions.  The same proof
chooses a prime $q>\deg P+r+2$ and filters the exponent
$q-2-r$, whose coefficient is
$(q-2)!/((q-2-r)!K'(1-q))\ne0$; for $P=0$, take any
$q>r+2$.

There is nevertheless a useful connection with the gap-series
literature.  Erd\H{o}s recorded a conjecture of Fej\'er and P\'olya
that every entire series $\sum_k a_kz^{n_k}$ with all $a_k\ne0$
and $n_k/k\to\infty$ attains every complex value infinitely
often \cite[Section IV.8, p.~250]{playErdos1961}.
For $f-m$ the exponents are $n_k=p_{k+1}-2\sim k\log k$.
The stronger summability condition $\sum_k1/n_k<\infty$ in the
proved Fej\'er-gap theorem does not hold here
\cite[Section 1]{playMurai1983}.
Our conclusion uses the extra residue-class structure of these
particular exponents; it makes no assertion about the general
Fej\'er--P\'olya conjecture.

Before finding this argument, we explored the middle height
numerically.  Truncate the Euler product defining $K$
at primes at most $Q$, call it $K_Q$, and use the odd polynomial
\[
 P_Q(z)=\sum_{3\le p\le251}\frac{z^{p-2}}{K_Q'(1-p)}.
\]
Numerical root finding gives the following nonzero roots in the
first quadrant.
\begin{center}
\begin{tabular}{rc}
\toprule
Euler cutoff $Q$ & Root of $P_Q$, rounded\tabularnewline
\midrule
$2\,000$  & $4.202432+3.410403\,i$\tabularnewline
$8\,000$  & $4.200425+3.410428\,i$\tabularnewline
$32\,000$ & $4.199995+3.410432\,i$\tabularnewline
\bottomrule
\end{tabular}
\end{center}
At each fixed Euler cutoff, replacing the Taylor cutoff $251$ by
$149$ leaves the displayed digits unchanged.  These are roots of
finite approximations, not certified roots of $f-m$; no bound on
the remaining Euler-product error is included.  The theorem now
proves infinitely many roots without locating them.  Their positions
and rate of accumulation are taken up in
Question~\ref{play:midpoint-geometry}.

The origin is the only real solution: a nonconstant analytic
function that is nonincreasing on the real line cannot take the
same value at two distinct real points.  Oddness of $f-m$ and its
real coefficients put every root off the axes into a quartet
$z,-z,\overline z,-\overline z$.  The infinitude theorem alone
does not determine how many roots belong to such quartets.

\subsection{A full-dimensional staircase that keeps an echo}

The primitive sampling law provides another small surprise.  Its
distribution function may be singular, strictly increasing and
supported on an interval.  Yet a frequency present in that measure
reappears unchanged at every doubling of its pitch.

\begin{proposition}[Dyadic echoes]\label{play:echoes}
Under the hypotheses of Theorem~\ref{samp:staircase}, the measure
$\nu_g$ is invariant under $T(y)=2y\pmod1$.  If any lane is
disabled, its Fourier transform does not tend to zero at infinity.
More precisely, with
\[
 \widehat\nu_g(t)=\int_0^1e^{2\pi ity}\,d\nu_g(y),
\]
there is an integer $k\ne0$ such that
$\widehat\nu_g(2^jk)=\widehat\nu_g(k)\ne0$ for every $j\ge0$.
\end{proposition}
\begin{proof}
In the sampling model \eqref{samp:model}, deleting the first
binary digit replaces $(U,V)$ by $(U+V,V)$ and shifts the sequence
of fair bits.  At every prime the map $(U,V)\mapsto(U+V,V)$
permutes the nonzero residue pairs; it preserves their uniform
distribution and their independence between primes.  The shifted
fair bits have the same independent law.  This proves invariance.
The binary expansion ambiguity has measure zero because $\nu_g$
is nonatomic.

Invariance gives $\widehat\nu_g(2k)=\widehat\nu_g(k)$ for
integer $k$.  If all coefficients with $k\ne0$ vanished, integration
against every trigonometric polynomial would agree with Lebesgue
measure.  Their density among continuous functions on the circle
would force $\nu_g=\mathrm{Leb}$.  This is impossible when a lane
is disabled, because \eqref{samp:decomposition} then has a
nonzero singular part.  Choose a nonzero coefficient and iterate.
\end{proof}

In particular, the explicit singular law from
\eqref{samp:explicit-mask} has Hausdorff dimension one but no
Fourier decay along an infinite sequence.  Its Fourier dimension
is zero.  The mechanism is the familiar one for a nonuniform
dilation-invariant measure; the sampling construction supplies a
specific arithmetic instance.

There is also more than one way for dimension to vary inside a
single sampling law.  As a small exact example, disable only the
lane of $3$.  Depending on the residues of the step $V$, the
proportion of free bits takes three values:
\begin{center}
\begin{tabular}{lcc}
\toprule
Condition on the primitive step & Probability & Free-bit proportion\tabularnewline
\midrule
$3\mid V$ & $1/4$ & $1$\tabularnewline
$3\nmid V,\ 2\nmid V$ & $1/2$ & $5/6$\tabularnewline
$3\nmid V,\ 2\mid V$ & $1/4$ & $2/3$\tabularnewline
\bottomrule
\end{tabular}
\end{center}
Indeed, $\Prob(p\mid V)=1/(p+1)$ independently at distinct
primes.  If $3\mid V$, the disabled lane is avoided.  If $V$ is
odd and not divisible by $3$, exactly one index in every six has
least prime factor $3$.  If $V$ is even and not divisible by $3$,
primitivity forces all progression values to be odd, and one index
in every three is closed.

For a fixed one of these periodic patterns, every compatible
$n$-digit cylinder has probability $2^{-dn+O(1)}$, where $d$ is
its free-bit proportion.  Distinct patterns having incompatible
forced-zero sets are almost surely eventually distinguished by
their fair bits.  In the finite mixture, compatible patterns with
larger $d$ contribute exponentially less.  Consequently the
dyadic local dimension of this law is $1$, $5/6$ or $2/3$ with
the respective probabilities in the table.  Here dyadic local
dimension means
\[
 d_\nu(y)=\lim_{n\to\infty}
      \frac{-\log_2\nu(I_n(y))}{n},
\]
when the limit exists, where $I_n(y)$ is the dyadic interval of
length $2^{-n}$ containing $y$.  This example shows why the
full-dimension statement in the sampling theorem does not assert
one local dimension almost everywhere.

The three values give an exact starting point for the infinite-mask
problem in Question~\ref{play:local-dimensions}.  The geometry already
varies inside one law, independently of which Goldbach answers occur.
